
% This LaTeX was auto-generated from MATLAB code.
% To make changes, update the MATLAB code and republish this document.

\documentclass{article}
\usepackage{graphicx}
\usepackage{color}

\sloppy
\definecolor{lightgray}{gray}{0.5}
\setlength{\parindent}{0pt}

\begin{document}

    
    
\subsection*{Contents}

\begin{itemize}
\setlength{\itemsep}{-1ex}
   \item Local Functions
\end{itemize}
\begin{verbatim}
function [v_fxn, policy_fxn, phi_dist, net_assets] = f_spe(z_grid, z_prob, u_fxn, u_prime_inv, u_prime, Params, a_next_grid, a_fine_grid)

% [v_fxn, policy_fxn, phi_dist, net_assets] =
%       f_spe(r, markov_chain, u_fxn, u_prime_inv, Params, a_grid, a_fine_grid)
%
% Computes the stationary partial equilibrium of the Aiyagari / Bewley model
% using the Endogenous Grid Method (EGM) and returns:
%
%   v_fxn      - (n_a_fine x n_z)  value function on fine grid
%   policy_fxn - (n_a_fine x n_z)  savings policy function on fine grid
%   phi_dist   - (n_a_fine x n_z)  stationary joint distribution phi(a,z)
%   net_assets - scalar, sum_i sum_j a_i * phi(a_i, z_j)  (net asset demand)
%
% Inputs:
%   r            - real interest rate
%   markov_chain - cell or struct: {z_grid (n_z x 1), z_prob (n_z x n_z)}
%   u_fxn        - utility function handle: u_fxn(c, Params)
%   u_prime_inv  - handle for (u')^{-1}: u_prime_inv(mu, Params)
%   Params       - struct with fields:
%                    .a_min, .a_max, .n_a, .curve  (grid)
%                    .gamma, .beta                 (preferences)
%                    .n_z                          (income states)
%                    .max_iter, .e_stop            (convergence)
%   a_grid       - (n_a x 1) coarse wealth grid  [not used internally;
%                  kept for interface compatibility]
%   a_fine_grid  - (n_a_fine x 1) fine grid for distribution / policy output

% ── Setup ─────────────────────────────────────────────────────────────────

R        = 1 + Params.r;

u_fxn_h       = @(c)  u_fxn(c, Params);
u_prime_inv_h = @(mu) u_prime_inv(mu, Params);
u_prime_h     = @(c) u_prime(c, Params);

% ── Initialisation ────────────────────────────────────────────────────────

a_1    = a_next_grid(1);

% Initial consumption: cash-on-hand minus minimum saving
c_init = zeros(Params.n_a, Params.n_z);
for iz = 1:Params.n_z
    c_init(:, iz) = max(R * a_next_grid + exp(z_grid(iz)) - a_1, 1e-6);
end

mu_grid = u_prime_h(c_init);

% Initial value function: perpetual consumption, discounted
V_grid = zeros(Params.n_a, Params.n_z);
for iz = 1:Params.n_z
    u0 = u_fxn_h(c_init(:, iz));
    V_grid(:, iz) = u0 / (1 - Params.beta);
end

policy_grid = repmat(a_next_grid, 1, Params.n_z);

error_val = Inf;
iteration = 0;

% ── EGM iteration ─────────────────────────────────────────────────────────
tic;
while error_val > Params.e_stop && iteration < Params.max_iter

    V_grid_old  = V_grid;
    mu_grid_old = mu_grid;

    %── Step 1: EGM backward step ─────────────────────────────────────────
    [a_current_grid, c_current_grid] = compute_EGM( ...
        Params, z_grid, z_prob, a_next_grid, mu_grid_old, u_prime_inv_h);

    %── Step 2: classify constrained / unconstrained points ───────────────
    constrained_mask = a_current_grid <= a_1;   % a_1 = a_next_grid(1)
    unconstrained_mask = ~constrained_mask;

    %── Step 3: value function update ─────────────────────────────────────
    V_grid_unconstrained = calc_unconstrained( ...
        Params, a_current_grid, V_grid_old, a_next_grid, unconstrained_mask);

    [V_grid_constrained, ~] = calc_constrained( ...
        Params, z_grid, z_prob, a_next_grid, V_grid_old, u_fxn_h);

    % Merge: constrained values override unconstrained interpolation
    V_grid = V_grid_unconstrained;
    V_grid(constrained_mask) = V_grid_constrained(constrained_mask);

    %── Step 4: update policy and marginal utility ────────────────────────
    c_new = zeros(Params.n_a, Params.n_z);

    for iz = 1:Params.n_z
        unc = unconstrained_mask(:, iz);
        con = constrained_mask(:, iz);

        % Unconstrained: spline-interpolate c from endogenous grid
        a_endog = a_current_grid(unc, iz);
        c_endog = c_current_grid(unc, iz);
        [a_s, si] = sort(real(a_endog));
        c_s       = c_endog(si);
        [a_u, ui] = unique(a_s);
        c_u       = c_s(ui);
        if length(a_u) >= 2
            c_new(unc, iz) = max( ...
                cubic_spline_interpolation(a_u, c_u, a_next_grid(unc)), 1e-10);
        else
            c_new(unc, iz) = max(c_endog, 1e-10);
        end

        % Constrained: c = R*a + exp(z) - a_1  (bind at borrowing limit)
        c_new(con, iz) = max(R * a_next_grid(con) + exp(z_grid(iz)) - a_1, 1e-10);

        % Policy: unconstrained saves a_current (endogenous), constrained saves a_1
        policy_grid(unc, iz) = a_next_grid(unc); % changed per ai
        policy_grid(con, iz) = a_1;
    end

    mu_grid = max(c_new, 1e-10) .^ (-Params.gamma);

    %── Step 5: convergence ────────────────────────────────────────────────
    error_val = max(abs(V_grid(:) - V_grid_old(:)));
    iteration = iteration + 1;
end
time_egm = toc;


fprintf('EGM converged after %d iterations in %.2f seconds (error = %.2e)\n', ...
        iteration, time_egm, error_val);

% ── Interpolate value and policy onto fine grid ───────────────────────────

n_a_fine   = length(a_fine_grid);

v_fxn      = zeros(n_a_fine, Params.n_z);
policy_fxn = zeros(n_a_fine, Params.n_z);

for iz = 1:Params.n_z
    v_fxn(:, iz)      = cubic_spline_interpolation(a_next_grid, V_grid(:, iz),      a_fine_grid);
    policy_fxn(:, iz) = cubic_spline_interpolation(a_next_grid, policy_grid(:, iz), a_fine_grid);
end

% ── Build state-transition matrix on fine grid ────────────────────────────

numstates = n_a_fine * Params.n_z;
policy_P  = zeros(numstates);

for iz = 1:Params.n_z
    for ia = 1:n_a_fine
        s_from  = (iz - 1) * n_a_fine + ia;

        a_prime = policy_fxn(ia, iz);
        a_prime = max(a_fine_grid(1), min(a_fine_grid(end), a_prime));

        % Index of the grid point immediately to the left of a_prime
        ik = min(find(a_fine_grid <= a_prime, 1, 'last'), n_a_fine - 1);

        % Linear interpolation weights
        w_lo = (a_fine_grid(ik+1) - a_prime)   / (a_fine_grid(ik+1) - a_fine_grid(ik));
        w_hi = (a_prime           - a_fine_grid(ik)) / (a_fine_grid(ik+1) - a_fine_grid(ik));

        for iz_next = 1:Params.n_z
            pi_z  = z_prob(iz, iz_next);
            s_lo  = (iz_next - 1) * n_a_fine + ik;
            s_hi  = (iz_next - 1) * n_a_fine + (ik + 1);
            policy_P(s_lo, s_from) = policy_P(s_lo, s_from) + pi_z * w_lo;
            policy_P(s_hi, s_from) = policy_P(s_hi, s_from) + pi_z * w_hi;
        end
    end
end

% ── Stationary distribution via eigenvalue method ─────────────────────────
% mc_invdist already issues a warning if the unit eigenvalue is not unique.

stationary_eigen = mc_invdist(policy_P');   % (numstates x 1), sums to 1

% Reshape to (n_a_fine x n_z)
phi_dist = reshape(stationary_eigen, n_a_fine, Params.n_z);

% ── Net asset demand ──────────────────────────────────────────────────────
% Net demand = sum_i sum_j  a_i * phi(a_i, z_j)
% = a_fine_grid' * (sum over z of phi_dist)

lambda_marginal_a = sum(phi_dist, 2);              % (n_a_fine x 1)
net_assets        = sum(a_fine_grid(:) .* lambda_marginal_a);


end
\end{verbatim}

        \color{lightgray} \begin{verbatim}Not enough input arguments.

Error in f_spe (line 30)
R        = 1 + Params.r;
               ^^^^^^^^\end{verbatim} \color{black}
    

\subsection*{Local Functions}

\begin{par}
-- cubic spline ──────────────────────────────────────────────────────
\end{par} \vspace{1em}
\begin{verbatim}
function V_interp = cubic_spline_interpolation(K_grid, V_grid, K_query)

% [V_interp = cubic_spline_interpolation(K_grid, V_grid, K_query)
%
%  Cubic spline interpolation between grid points (for value function)
%
% K_grid: existing / known grid points (e.g. asset grid)
% V_grid: value function evaluated at K_grid points / or function to be eval at
% K_query: new grid points where we want to evaluate the value function
%
% reference: https://www.mathworks.com/help/matlab/ref/spline.html

    V_grid_clean = V_grid;
    V_grid_clean(~isfinite(V_grid_clean)) = -1e10;

    spline_interp = spline(K_grid, V_grid_clean);
    V_interp = ppval(spline_interp, K_query);
end


% -- compute EGM ──────────────────────────────────────────────────────
function [a_current_grid, c_current_grid] = compute_EGM(Params, z_grid, z_prob, a_next_grid, mu_grid, u_prime_inv)

% [a_current_grid, c_current_grid] = compute_EGM(Params, z_grid, z_prob,
%                                                  a_next_grid, mu_grid, u_prime_inv)
%
% EGM backward step. Given the exogenous grid a' and the current marginal
% utility grid mu(a',z) = u'(c(a',z)) = c(a',z)^{-gamma}, recovers the
% endogenous current-period wealth grid a(a',z) by inverting the Euler eq:
%
%   Euler:  u'(c) = beta*(1+r)*E[u'(c(a',z'))|z]
%   =>  c(a',z) = [beta*(1+r)*E[mu(a',z')|z]]^{-1/gamma}
%   =>  a(a',z) = (c + a' - exp(z)) / (1+r)        [budget constraint]
% ───────────────────────────────────────────────────────────────────────────
%
% Inputs:
%   Params      - parameter struct (needs .beta, .r, .n_a, .n_z)
%   z_grid      - (n_z x 1) discretised log-income grid
%   z_prob      - (n_z x n_z) transition matrix, rows sum to 1
%   a_next_grid - (n_a x 1) exogenous savings grid
%   mu_grid     - (n_a x n_z) marginal utility u'(c(a',z)) on exog grid
%   u_prime_inv - function handle: inverse of u', i.e. (u')^{-1}
%
% Outputs:
%   a_current_grid - (n_a x n_z) endogenous current-wealth grid
%   c_current_grid - (n_a x n_z) consumption implied by Euler equation

    beta = Params.beta;
    R    = 1 + Params.r;

    % EMU(ia,iz) = sum_{iz'} pi(iz,iz') * mu(ia,iz')
    % mu_grid is (n_a x n_z), z_prob is (n_z x n_z)
    % EMU = mu_grid * z_prob'   gives (n_a x n_z) — correct orientation
    EMU = mu_grid * z_prob';

    a_current_grid = zeros(Params.n_a, Params.n_z);
    c_current_grid = zeros(Params.n_a, Params.n_z);

    for ia = 1:Params.n_a
        a_next = a_next_grid(ia);
        for iz = 1:Params.n_z
            z_curr = z_grid(iz);


            % Euler: u'(c) = beta*R * E[mu]  =>  c = (u')^{-1}(beta*R * E[mu])
            c_curr = max(u_prime_inv(beta * R * EMU(ia, iz)), 1e-10);

            % Budget constraint: c + a' = R*a + exp(z)  =>  a = (c + a' - exp(z)) / R
            a_curr = (a_next - exp(z_curr) + c_curr) / R;

            a_current_grid(ia, iz) = a_curr;
            c_current_grid(ia, iz) = c_curr;
        end
    end
end

% -- calc_contrained ──────────────────────────────────────────────────────

function [V_grid_constrained, policy_grid_constrained] = calc_constrained( ...
            Params, z_grid, z_prob, a_next_grid, V_grid_old, u_fxn)

% [V_grid_constrained, policy_grid_constrained] = calc_constrained(...)
%
% Evaluates the Bellman equation for ALL current-wealth levels a_i under
% the CONSTRAINED policy a' = a_next_grid(1) = a_min (borrowing constraint).
%
% Used for grid points where the unconstrained EGM solution gives
% a_current <= a_min, i.e. the household would like to borrow more than
% the constraint permits.
%
% Formula:
%   V(a_i, z_j) = u(R*a_i + exp(z_j) - a_1) + beta * E[V(a_1, z') | z_j]
%
% where a_i are the CURRENT asset holdings (= a_next_grid used as the
% current-wealth vector in this evaluation).
%

%
% ── Inputs ──────────────────────────────────────────────────────────────────
%   Params      - parameter struct (.r, .beta, .n_a, .n_z)
%   z_grid      - (n_z x 1) log-income grid
%   z_prob      - (n_z x n_z) transition matrix
%   a_next_grid - (n_a x 1) exogenous savings grid; a_1 = a_next_grid(1)
%   V_grid_old  - (n_a x n_z) value function on exogenous grid (previous iter)
%   u_fxn       - function handle: u_fxn(c, Params)

    a_1 = a_next_grid(1);   % constrained saving = grid minimum
    R   = 1 + Params.r;

    % Continuation value at a_1 for each current income state:
    %   EV_a1(iz) = sum_{iz'} pi(iz, iz') * V(a_1, iz')
    % V_grid_old(1,:) is V evaluated at a' = a_1 for every income state iz'.
    EV_a1 = z_prob * V_grid_old(1, :)';   % (n_z x 1)

    V_grid_constrained      = zeros(Params.n_a, Params.n_z);
    policy_grid_constrained = ones(Params.n_a,  Params.n_z) * a_1;

    for iz = 1:Params.n_z
        z_curr = z_grid(iz);

        % Consumption at each current-wealth level a_i when a' is fixed at a_1:
        %   c_i = R * a_i + exp(z) - a_1
        % Here a_next_grid stands in for the current-wealth vector a_i
        % (the exogenous grid is used for both roles in EGM).
        c = R * a_next_grid + exp(z_curr) - a_1;   % (n_a x 1)

        u = u_fxn(c);

        % EV_a1(iz) is a scalar — the same continuation value for all ia
        % because every constrained household saves exactly a_1.
        V_grid_constrained(:, iz) = u + Params.beta * EV_a1(iz);
    end
end

% -- calc_uncontrained ──────────────────────────────────────────────────────
function [V_grid_unconstrained] = calc_unconstrained(Params, a_current_grid, V_grid_old, a_next_grid, unconstrained_mask)

% [V_grid_unconstrained] = calc_unconstrained(...)
%
% Unconstrained EGM update: interpolates V from the endogenous grid
% a_current_grid back onto the fixed exogenous grid a_next_grid.
%
% For each income state iz, the endogenous grid points a_current_grid(:,iz)
% are the x-knots and V_grid_old(:,iz) are the corresponding values.
% We fit a spline and evaluate at the fixed a_next_grid query points.
%
% Inputs:
%   a_current_grid   - (n_a x n_z) endogenous current-wealth grid from EGM
%   V_grid_old       - (n_a x n_z) value function on exogenous grid
%   a_next_grid      - (n_a x 1)   fixed exogenous savings grid
%   unconstrained_mask - (n_a x n_z) logical, true where a_current > a_min

    V_grid_unconstrained = zeros(Params.n_a, Params.n_z);

    for iz = 1:Params.n_z
        unc_rows = unconstrained_mask(:, iz);

        a_endog = a_current_grid(unc_rows, iz);
        V_endog = V_grid_old(unc_rows, iz);

        [a_sorted, sort_idx] = sort(a_endog);
        V_sorted             = V_endog(sort_idx);

        [a_unique, unique_idx] = unique(a_sorted);
        V_unique               = V_sorted(unique_idx);

        a_unique = real(a_unique);
        V_unique = real(V_unique);

        % need at least 2 points to interpolate
        if length(a_unique) < 2
            V_grid_unconstrained(:, iz) = V_grid_old(:, iz);
            continue
        end

        % interpolate on full grid
        V_grid_unconstrained(:, iz) = cubic_spline_interpolation(a_unique, V_unique, a_next_grid);
    end
end

% -- mc_invdist ──────────────────────────────────────────────────────
function P = mc_invdist(PI)

% P = mc_invdist(PI)
%
% Computes invariant distribution P of a Markov chain with transition
% matrix PI

[V, D] = eig(PI');

ii = find(abs(diag(D) - 1) < 1E-8, 1);      % Find first unit eigenvalue

P = V(:,ii) / sum(V(:,ii));                 % Normalize unit eigenvector

assert(max(abs(P' - P'*PI)) < 1E-12)        % Verify dist. is stationary

if sum(abs(diag(D) - 1) < 1E-8) > 1         % Check for uniqueness
   warning('P not unique')
end

end
\end{verbatim}



\end{document}

